%% Branching of the Hydrogen Cascade after Lyman-Series Absorption
%% Author: Kwang-Il Seon
\documentclass[english,a4paper]{my_memo}
\synctex=1
\usepackage{amsmath}
\usepackage{amssymb}
\usepackage{microtype}
\usepackage{babel}
\usepackage{booktabs}

\begin{document}

\title{Branching of the Hydrogen Cascade after Lyman-Series Absorption}
\author{Kwang-Il Seon}
\date{\docmoddate}
\maketitle

\tableofcontents
\bigskip

%=============================================================================
\section{Introduction and Assumptions}
%=============================================================================

A Lyman-series absorption event promotes the hydrogen electron from the
ground state to an excited $p$ level,
\begin{equation}
   1s \;\longrightarrow\; np , \qquad n \ge 2 ,
\end{equation}
which corresponds to Ly$\alpha$ ($n=2$), Ly$\beta$ ($n=3$), Ly$\gamma$
($n=4$), and so on.  Once in the $np$ level, the electron decays
radiatively.  The allowed electric-dipole downward channels from $np$ are
to the $ms$ states ($m<n$) and to the $md$ states ($3 \le m < n$).  Each
downward step lands the electron in a new bound level, from which it decays
again; the process repeats until the electron reaches the ground state.  We
call the complete sequence of radiative transitions from the initially
excited $np$ level down to $1s$ the \emph{cascade}.

Two terminal channels close the cascade and are physically distinct:
\begin{itemize}
   \item[(a)] \textbf{Direct Lyman emission.}  Whenever the electron reaches
     an $n'p$ level, it may decay directly to the ground state,
     $n'p \to 1s$, emitting a Lyman-series photon.  This ends the cascade on
     the ground state.
   \item[(b)] \textbf{Two-photon continuum.}  The cascade may instead route
     the electron into the metastable $2s$ state.  Because $2s \to 1s$ is
     electric-dipole forbidden, this level decays by \emph{two-photon}
     emission, $A(2s \to 1s) = 8.229~\mathrm{s}^{-1}$, producing a broad
     continuum rather than a spectral line.
\end{itemize}

Throughout we adopt the low-density nebular (Case-B-like) regime: a pure
radiative cascade with no collisional $l$-redistribution, and with the $2s$
state always decaying via two-photon emission.  In this limit the branching
of the cascade is fixed entirely by the spontaneous transition
probabilities of hydrogen, independent of density.  A useful example is
Ly$\beta$: absorption into $3p$ can decay directly back to $1s$
(re-emitting Ly$\beta$) or route through $3p \to 2s$, which is a component
of H$\alpha$.  In the latter case the Ly$\beta$ absorption is physically
converted into an H$\alpha$ photon ($3p \to 2s$) plus a two-photon
($2s \to 1s$) continuum.

%=============================================================================
\section{Lyman-Series Oscillator Strengths}
%=============================================================================

Table~\ref{tab:fosc} lists the absorption oscillator strengths
$f(1s \to np)$ of the hydrogen Lyman series, together with the transition
wavelengths and the ratios relative to Ly$\alpha$.  For large $n$ the
oscillator strength falls off asymptotically as $n^{-3}$, so successive
Lyman lines weaken rapidly.

It is important to note that the line-center opacity scales as the product
$f\,\lambda$, not as $f$ alone.  Because the higher Lyman lines lie at
shorter wavelengths, the opacity ratio is slightly smaller than the
oscillator-strength ratio.  For Ly$\beta$ relative to Ly$\alpha$,
\begin{equation}
   \frac{\tau_0(\mathrm{Ly}\beta)}{\tau_0(\mathrm{Ly}\alpha)}
   = 0.1901 \times \frac{1025.72}{1215.67}
   = 0.160 ,
\end{equation}
rather than $0.190$.

\begin{table}[htb]
\centering
\caption{Lyman-series absorption oscillator strengths $f(1s \to np)$.}
\label{tab:fosc}
\small
\begin{tabular}{c l c c c}
\toprule
$n$ & line & $\lambda$ [\AA] & $f$ & $f/f(\mathrm{Ly}\alpha)$ \\
\midrule
2  & Ly$\alpha$      & 1215.67 & 0.4162   & 1.00000  \\
3  & Ly$\beta$       & 1025.72 & 0.07910  & 0.19006  \\
4  & Ly$\gamma$      & 972.54  & 0.02899  & 0.06966  \\
5  & Ly$\delta$      & 949.74  & 0.01394  & 0.03349  \\
6  & Ly$\epsilon$    & 937.80  & 0.007799 & 0.01874  \\
7  & $1 \to 7$       & 930.75  & 0.004814 & 0.01157  \\
8  & $1 \to 8$       & 926.23  & 0.003183 & 0.00765  \\
9  & $1 \to 9$       & 923.15  & 0.002216 & 0.00532  \\
10 & $1 \to 10$      & 920.96  & 0.001605 & 0.00386  \\
\bottomrule
\end{tabular}
\end{table}

%=============================================================================
\section{Method and Validation}
%=============================================================================

The spontaneous transition probabilities $A(nl \to n'l')$ were computed from
the exact hydrogenic dipole radial integrals and calibrated on the
analytically known value $A(2p \to 1s) = 6.2649 \times 10^{8}~\mathrm{s}^{-1}$,
which fixes the Einstein-$A$ prefactor.  The absorption oscillator strengths
follow from
\begin{equation}
   f(1s \to np) = \tfrac{2}{3}\,\Delta E\;\big|\langle np|\,r\,|1s\rangle\big|^2 ,
\end{equation}
evaluated in atomic units, where $\Delta E$ is the transition energy.  Every
computed value reproduces the literature to four or five significant
figures.  Table~\ref{tab:validate} demonstrates this agreement for a
representative set of transitions spanning the Lyman, Balmer, and Paschen
series; the computed and literature $A$ coefficients agree to within
$0.0$--$0.4\%$.

\begin{table}[htb]
\centering
\caption{Validation of computed Einstein $A$ coefficients (s$^{-1}$).}
\label{tab:validate}
\small
\begin{tabular}{l l r r}
\toprule
transition & series & computed & literature \\
\midrule
$A(2p \to 1s)$ & Ly$\alpha$  & $6.2649 \times 10^{8}$ & $6.2649 \times 10^{8}$ \\
$A(3p \to 1s)$ & Ly$\beta$   & $1.6725 \times 10^{8}$ & $1.6725 \times 10^{8}$ \\
$A(3s \to 2p)$ & H$\alpha$   & $6.3136 \times 10^{6}$ & $6.3132 \times 10^{6}$ \\
$A(3p \to 2s)$ & H$\alpha$   & $2.2448 \times 10^{7}$ & $2.2449 \times 10^{7}$ \\
$A(3d \to 2p)$ & H$\alpha$   & $6.4651 \times 10^{7}$ & $6.4651 \times 10^{7}$ \\
$A(4p \to 1s)$ & Ly$\gamma$  & $6.8187 \times 10^{7}$ & $6.8186 \times 10^{7}$ \\
$A(4d \to 2p)$ & H$\beta$    & $2.0625 \times 10^{7}$ & $2.0625 \times 10^{7}$ \\
$A(4f \to 3d)$ & Pa$\alpha$  & $1.3788 \times 10^{7}$ & $1.3785 \times 10^{7}$ \\
\bottomrule
\end{tabular}
\end{table}

%=============================================================================
\section{Direct Branching Ratios of the $np$ Levels}
%=============================================================================

Table~\ref{tab:branch} gives the direct branching ratios
$b(np \to n'l') = A(np \to n'l')/A_{\rm tot}(np)$ of the initially excited
$np$ level, expressed as percentages.  Here $A_{\rm tot}$ is the total
spontaneous decay rate and $\tau = 1/A_{\rm tot}$ the radiative lifetime.
These are the single-step probabilities out of the $np$ level, not yet the
full cascade outcome (Section~\ref{sec:cascade}).

Every $np$ level has exactly three kinds of downward channel, and the table
groups them so that the three add up to $100\,\%$:
(i)~the resonance line $np \to 1s$, a Lyman photon straight back to the ground
state;
(ii)~$np \to 2s$, which populates the metastable $2s$ level and therefore
feeds the two-photon continuum;
(iii)~the subordinate decays to any lower level with $m \ge 3$ (the $ms$ and
$md$ states $3s, 3d, 4s, 4d, \dots$), from which the cascade continues.
The last two columns single out the two largest subordinate channels,
$np \to 3s$ and $np \to 3d$; these are \emph{already included} in the
$m \ge 3$ column and are listed only to show where the subordinate weight
goes.

From $3p$ only the first two channels exist, because no $md$ state with
$m \ge 3$ lies below $n = 3$.  Subordinate decay first appears at $4p$ and
grows with $n$: the Lyman branch $np \to 1s$ falls monotonically from
$100\,\%$ at $2p$ to $78.5\,\%$ at $10p$, the subordinate fraction rises to
about $10\,\%$, and the two-photon feed $np \to 2s$ stays nearly constant
near $11.5$--$11.9\,\%$ across the whole series.

\begin{table}[htb]
\centering
\caption{Direct (single-step) decay of the $np$ levels, as branching
percentages.  $\tau = 1/A_{\rm tot}$ is the radiative lifetime.  The three
columns under ``Branching ratio'' are the probabilities that the first
emitted photon is the Lyman line $np \to 1s$, feeds the metastable $2s$
(and hence the two-photon continuum), or goes to a subordinate level with
lower principal quantum number $m \ge 3$; \emph{these three sum to
$100\,\%$}.  The two right-hand columns break out the two largest
subordinate channels, $np \to 3s$ and $np \to 3d$, which are already counted
in the $m \ge 3$ column.  ``--'' marks a channel that does not exist.}
\label{tab:branch}
\footnotesize
\setlength{\tabcolsep}{5pt}
\begin{tabular}{l r r c c c c c}
\toprule
      &            &                      & \multicolumn{3}{c}{Branching ratio (\%)} & \multicolumn{2}{c}{Largest subord. (\%)} \\
\cmidrule(lr){4-6}\cmidrule(lr){7-8}
$np$  & $\tau$ [ns] & $A_{\rm tot}$ [s$^{-1}$] & $\to 1s$ & $\to 2s$ & $\to m\!\ge\!3$ & $\to 3s$ & $\to 3d$ \\
      &            &                      & (Lyman)  & (2-photon) & (subord.)     &          &          \\
\midrule
$2p$  & 1.596   & $6.2649 \times 10^{8}$ & 100.00 & --    & --    & --   & --   \\
$3p$  & 5.271   & $1.8970 \times 10^{8}$ & 88.17  & 11.83 & --    & --   & --   \\
$4p$  & 12.305  & $8.1267 \times 10^{7}$ & 83.90  & 11.90 & 4.20  & 3.77 & 0.43 \\
$5p$  & 23.789  & $4.2037 \times 10^{7}$ & 81.78  & 11.77 & 6.45  & 3.90 & 0.36 \\
$6p$  & 40.819  & $2.4499 \times 10^{7}$ & 80.53  & 11.67 & 7.80  & 3.90 & 0.32 \\
$7p$  & 64.492  & $1.5506 \times 10^{7}$ & 79.72  & 11.59 & 8.69  & 3.89 & 0.30 \\
$8p$  & 95.909  & $1.0427 \times 10^{7}$ & 79.17  & 11.53 & 9.30  & 3.87 & 0.29 \\
$9p$  & 136.169 & $7.3438 \times 10^{6}$ & 78.77  & 11.49 & 9.74  & 3.86 & 0.28 \\
$10p$ & 186.372 & $5.3656 \times 10^{6}$ & 78.47  & 11.46 & 10.07 & 3.85 & 0.27 \\
\bottomrule
\end{tabular}
\end{table}

%=============================================================================
\section{Cascade Outcome Probabilities}
\label{sec:cascade}
%=============================================================================

Following the cascade all the way to the ground state, every absorbed
Lyman-series photon ends in exactly one of the two terminal channels of
Section~1: either the $2s$ two-photon continuum, or a direct Lyman photon
$n'p \to 1s$.  Table~\ref{tab:outcome} gives these total probabilities as a
function of the initially excited level.

The most striking feature is the near-constancy of the outcome.  The
probability of ending in the two-photon continuum is
\begin{equation}
   P(\text{two-photon}) \approx 0.118 \quad (11.8\%) ,
\end{equation}
essentially independent of the initial principal quantum number $n$, and
correspondingly
\begin{equation}
   P(\text{direct Lyman}) \approx 0.88 .
\end{equation}
Ly$\alpha$ ($2p$) is the sole exception: $2p$ can only decay to $1s$, so it
never routes through $2s$ and always re-emits a Lyman photon.  For all
higher Lyman lines, roughly one absorbed photon in eight is converted into
the two-photon continuum, with the remainder re-emerging as a direct Lyman
photon.  The two-photon probability actually peaks at $4p$ ($0.11897$) and
then stays remarkably flat, remaining within $0.1180$--$0.1190$ across the
entire series through $10p$; the direct-Lyman probability correspondingly
sits at $\approx 0.881$--$0.882$.

\begin{table}[htb]
\centering
\caption{Cascade outcome after $1s \to np$ absorption: total probability of
the $2s$ two-photon continuum versus a direct Lyman photon.}
\label{tab:outcome}
\small
\begin{tabular}{l c c}
\toprule
initial level & $P(\text{2-photon via } 2s)$ & $P(\text{direct Lyman } n'p \to 1s)$ \\
\midrule
$2p$  (Ly$\alpha$)   & 0.00000 & 1.00000 \\
$3p$  (Ly$\beta$)    & 0.11834 & 0.88166 \\
$4p$  (Ly$\gamma$)   & 0.11897 & 0.88103 \\
$5p$  (Ly$\delta$)   & 0.11871 & 0.88129 \\
$6p$  (Ly$\epsilon$) & 0.11848 & 0.88152 \\
$7p$  ($1 \to 7$)    & 0.11831 & 0.88169 \\
$8p$  ($1 \to 8$)    & 0.11818 & 0.88182 \\
$9p$  ($1 \to 9$)    & 0.11808 & 0.88192 \\
$10p$ ($1 \to 10$)   & 0.11800 & 0.88200 \\
\bottomrule
\end{tabular}
\end{table}

%=============================================================================
\section{Cascade Route Decomposition}
\label{sec:routes}
%=============================================================================

Rather than counting how often each individual transition is emitted, it is
far more illuminating to follow the complete \emph{routes} the electron takes
from the initially excited $np$ level down to the ground state.  A route is a
chain $np \to \cdots \to 1s$; the transitions along it are exactly the photons
emitted \emph{together} in one de-excitation, and the route probability is the
product of the branching ratios of its steps.  Because every absorbed photon
follows one and only one route to the ground state, the route probabilities
sum to unity.

The chains are forced by the dipole selection rule $\Delta l = \pm 1$: only
$p$ states connect to the $1s$ ground state, so an electron that reaches an
$s$ or $d$ state must first step to a lower $p$ state before it can reach
$1s$.  This is why $3s$ does not decay to $1s$ directly but must follow
$3s \to 2p \to 1s$, and it is what makes the routes fall into a few
systematic families.

Table~\ref{tab:routes} lists the dominant routes, organized into five
families:
(A)~the direct resonance decay $np \to 1s$;
(B)~the metastable route $np \to 2s \to 1s$, which ends in the two-photon
continuum;
(C)~routes through one subordinate $ns$ state that end on Ly$\alpha$
($\cdots \to 2p \to 1s$);
(D)~routes through one subordinate $nd$ state that end on Ly$\alpha$; and
(E)~routes through an $ns$ state that end on Ly$\beta$
($\cdots \to 3p \to 1s$).
The intermediate steps make the co-emitted photons explicit: e.g.\ the route
$np \to 3s \to 2p \to 1s$ emits, in succession, a subordinate photon, an
H$\alpha$ photon ($3s \to 2p$), and finally a Ly$\alpha$ photon
($2p \to 1s$).  All routes not listed are collected in the ``other routes''
residual; each has a probability below $\sim 0.3\%$, and they proliferate with
$n$ (the $10p$ cascade has $1157$ distinct routes in all).

Reading off the $10p$ column, for example: the absorbed Ly-10 photon is simply
re-emitted as $10p \to 1s$ with probability $78.5\%$; it is converted into the
two-photon continuum via $10p \to 2s \to 1s$ with probability $11.5\%$; and it
degrades into an H$\alpha$-then-Ly$\alpha$ pair via $10p \to 3s \to 2p \to 1s$
with probability $3.85\%$.  The route probabilities are remarkably insensitive
to the initial level --- each family contributes nearly the same fraction
whether the absorbed line is Ly$\gamma$ or Ly-10.

\begin{table}[htb]
\centering
\caption{Dominant cascade routes after $1s \to np$ absorption.  Each row is a
route --- a chain $np \to \cdots \to 1s$ whose transitions are the photons
emitted \emph{together} in one de-excitation --- and the entries are the route
probability (in \%) for each initially excited level.  Routes are grouped into
families (A)--(E) by their structure; ``$1s\,(2\gamma)$'' marks the
$2s \to 1s$ two-photon continuum.  ``--'' means the route is not open at that
$n$.  Every column, including the ``other routes'' residual, sums to
$100\%$.}
\label{tab:routes}
\footnotesize
\setlength{\tabcolsep}{4pt}
\begin{tabular}{l r r r r r r r r}
\toprule
route & $3p$ & $4p$ & $5p$ & $6p$ & $7p$ & $8p$ & $9p$ & $10p$ \\
\midrule
\multicolumn{9}{l}{\textit{(A) Direct resonance decay (final photon: the Ly-$n$ line)}}\\
$np \to 1s$                    & 88.17 & 83.90 & 81.78 & 80.53 & 79.72 & 79.17 & 78.77 & 78.47 \\
\addlinespace[2pt]
\multicolumn{9}{l}{\textit{(B) Via metastable $2s$ (two-photon continuum)}}\\
$np \to 2s \to 1s\,(2\gamma)$  & 11.83 & 11.90 & 11.77 & 11.67 & 11.59 & 11.53 & 11.49 & 11.46 \\
\addlinespace[2pt]
\multicolumn{9}{l}{\textit{(C) Via $ns$, ending on Ly$\alpha$ $(\to 2p \to 1s)$}}\\
$np \to 3s \to 2p \to 1s$      & --    & 3.77  & 3.90  & 3.90  & 3.89  & 3.87  & 3.86  & 3.85  \\
$np \to 4s \to 2p \to 1s$      & --    & --    & 1.02  & 1.06  & 1.07  & 1.06  & 1.06  & 1.06  \\
$np \to 5s \to 2p \to 1s$      & --    & --    & --    & 0.45  & 0.47  & 0.47  & 0.46  & 0.46  \\
$np \to 6s \to 2p \to 1s$      & --    & --    & --    & --    & 0.25  & 0.26  & 0.26  & 0.25  \\
\addlinespace[2pt]
\multicolumn{9}{l}{\textit{(D) Via $nd$, ending on Ly$\alpha$ $(\to 2p \to 1s)$}}\\
$np \to 3d \to 2p \to 1s$      & --    & 0.43  & 0.36  & 0.32  & 0.30  & 0.29  & 0.28  & 0.27  \\
$np \to 4d \to 2p \to 1s$      & --    & --    & 0.33  & 0.29  & 0.26  & 0.24  & 0.23  & 0.23  \\
$np \to 5d \to 2p \to 1s$      & --    & --    & --    & 0.26  & 0.22  & 0.20  & 0.19  & 0.18  \\
\addlinespace[2pt]
\multicolumn{9}{l}{\textit{(E) Via $ns$, ending on Ly$\beta$ $(\to 3p \to 1s)$}}\\
$np \to 4s \to 3p \to 1s$      & --    & --    & 0.64  & 0.67  & 0.67  & 0.67  & 0.67  & 0.66  \\
$np \to 5s \to 3p \to 1s$      & --    & --    & --    & 0.28  & 0.29  & 0.29  & 0.29  & 0.29  \\
\addlinespace[2pt]
\multicolumn{9}{l}{\textit{Other routes (each below $0.3\%$)}}\\
other routes                   & 0.00  & 0.00  & 0.20  & 0.58  & 1.28  & 1.95  & 2.44  & 2.81  \\
\bottomrule
\end{tabular}
\end{table}

%=============================================================================
\section{Discussion}
%=============================================================================

The near-constant $11.8\%$ conversion probability has a direct physical
interpretation.  For Ly$\beta$, an absorption that does not simply re-emit
Ly$\beta$ is instead converted, with probability $0.11834$, into an
H$\alpha$ photon ($3p \to 2s$) plus a $2s \to 1s$ two-photon continuum.  The
higher Lyman lines behave almost identically: about one absorbed photon in
eight is permanently removed from the Lyman series and redistributed into
subordinate lines and the two-photon continuum, while the remaining
$\sim\!88\%$ re-emerges as a direct Lyman photon.

Extending the analysis across the whole Lyman series through Ly-10
($1 \to 10$) reinforces this picture.  The conversion probability to the $2s$
two-photon continuum stays near $11.8$--$11.9\%$, essentially independent of
which Lyman line is absorbed.  What does change with $n$ is the fraction of
the absorbed photon's energy that is degraded into subordinate (Balmer,
Paschen, \dots) lines during the longer cascade: this grows steadily with the
initial principal quantum number, so that the total number of emitted photons
per absorbed Lyman-series photon rises from exactly $1$ at Ly$\alpha$ to
$\approx 1.20$ at Ly-10.  A higher Lyman-line absorption therefore returns a
direct Lyman photon with nearly the same probability regardless of $n$, but it
distributes progressively more of the cascade energy into the lower series as
$n$ increases.

This branching is central to the treatment of nebular and intergalactic
Lyman pumping, and to Ly$\alpha$ radiative transfer more generally.  Higher
Lyman-series photons that are absorbed are not conservatively scattered:
a fixed fraction is destroyed as line photons and converted into the
two-photon continuum and lower-series lines (for example, feeding H$\alpha$
and the two-photon continuum from Ly$\beta$).  In a Monte Carlo transfer
calculation this means that the excitation of $np$ levels by higher Lyman
lines cannot be modeled as pure resonance scattering; the cascade branching
of Tables~\ref{tab:outcome} and~\ref{tab:routes} determines how much Lyman
flux is conserved and how much is converted, essentially independently of the
initial level for $n \ge 3$.

\end{document}
